\section*{The Propers: Detailed Commentary}

\subsection*{The Gospel's controversy: mercy and false justice\enspace\properrefs{Gosp., Comm.}}

Luke first names the conflict: publicans and sinners draw near to hear Jesus, while Pharisees and scribes murmur because he receives and eats with them. The parables answer an accusation about Christ's table-fellowship. They therefore reveal both the seeker's character and the hearer's possible resistance.

\begin{fourstable}{Witness}{Direct locus}{Exegetical emphasis}{Precise sense}
St. Cyril of Alexandria & \emph{Commentary on Luke}, Homily 106, on Luke 15:1--10 & Christ receives those who approach because he is the saving physician. The accusers make mercy itself an accusation; Christ answers by requiring them to share the joy caused by repentance. & Cyril does not turn welcome into approval of sin. The movement is approach, healing, repentance, and joy.\\
St. Gregory the Great & \emph{Forty Gospel Homilies}, Homily 34 & True justice pities; false justice disdains. Converted sinners may love with unusual ardor because memory of their wandering intensifies gratitude. & Gregory expressly avoids making every penitent superior to every person who remained just. His contrast diagnoses dispositions, not a mechanical ranking.\\
Tertullian & \emph{On Patience}, 12, receiving Luke 15 & Divine patience goes after one sheep and bears it; impatience would have despised the one for the sake of the many. Patience is active pursuit, not mere delay. & Tertullian receives the parable early and ascetically; he comments neither on the Alleluia nor on this Roman formulary.\\
\end{fourstable}

\begin{sourcecard}{At the contested table}
Cyril begins with the sinners' willingness to hear and treats Christ's nearness as medicinal: the Savior does not participate in their disorder but makes table-fellowship an occasion of rescue. Gregory addresses the opposite disposition, in which a reputation for exactness hardens into contempt. Read together, the witnesses distinguish welcome ordered to healing from both moral indifference and self-protective disdain.
\end{sourcecard}

\begin{studybox}{Mercy has a definite shape}
The Gospel does not oppose justice to indifference. The Shepherd finds; the sinner repents; heaven rejoices. Mercy's initiative makes conversion possible, and conversion discloses rather than purchases mercy. The Pharisaic error is not zeal for holiness as such, but refusal of the holy joy proper to a sinner's restoration.
\end{studybox}

\sectionguard
\subsection*{The Shepherd, the shoulders, and the home\enspace\properrefs{Gosp., Comm.}}

The first parable is deliberately disproportionate: one is missing, and the seeker refuses to treat that loss as negligible. Patristic interpretation enlarges the one without erasing Luke's direct appeal to each sinner.

\begin{threestable}{Patristic movement}{Textual ground}{Theological contribution}
One sheep as fallen humanity & Gregory, Homily 34, reads the hundred as the fullness of rational creation, the ninety-nine as the angelic hosts remaining above, and the one as humanity wandering through sin. & The Shepherd's descent is not a minor adjustment to an otherwise complete world. In the one human nature, the whole race is sought.\\
The shoulders as assumed nature and Cross & Gregory connects the carrying to Christ's assumption of human nature and bearing of sin; Ambrose, \emph{Exposition of Luke} VII.208--210, makes the shoulders the arms of the Cross. & Recovery rests on Christ's labor. The sheep does not climb back by autonomous strength; the Incarnate and Crucified bears it home.\\
One body in Adam and Christ & Ambrose can speak of the single sheep corporately: humanity was lost in Adam and is gathered in Christ. & Personal repentance belongs inside a common human wound and a common ecclesial restoration. Individual concern and corporate salvation are not rivals.\\
Friends and neighbors & Gregory identifies the rejoicing circle with the angelic hosts; Ambrose likewise places restored humanity amid heavenly joy. & The Communion antiphon is not decorative sentiment. Eucharistic assembly rehearses the communion into which the sinner is carried.\\
\end{threestable}

\begin{studybox}{The hundred, the ninety-nine, and the one}
The Fathers' hundred and ninety-nine are theological figures of fullness and angelic communion, not a census of saved and lost persons. Nor does one sinner make salvation private: Luke's singular protects the irreducible value of the person, while Gregory and Ambrose show the one human nature borne home in Christ. Ambrose's larger reading of the sheep, coin, and son as a threefold remedy associates carrying, searching, and receiving with one saving work rather than three rival agents.
\end{studybox}

\sectionguard
\subsection*{The drachma: three searches within the house\enspace\properrefs{Gosp., Comm.}}

The second parable repeats the sequence---loss, initiative, finding, convoked joy---but changes the field. The coin is not roaming outside; it is hidden within a house and bears an image whose value remains even when obscured. Gregory, Ambrose, and Gregory of Nyssa use those features differently.

\begin{fourstable}{Witness}{Woman and lamp}{Coin and house}{Distinctive result}
St. Gregory the Great, Homily 34 & The woman signifies divine Wisdom; the lamp is the manifestation of divinity in the earthen vessel of Christ's flesh. Sweeping names the searching disturbance by which conscience is exposed. & Fallen humanity lies hidden in the house of the world, bearing the Creator's stamp but obscured by sin. & Incarnation illuminates; searching grace makes the hidden image visible and restores it to heavenly fellowship.\\
St. Ambrose, \emph{Luke} VII.211--212 & The woman can signify the Church actively searching with the lamp. & The drachma bears the sovereign's image and belongs to the Church's treasure. Human dignity is received, not self-minted. & Ecclesial proclamation seeks what belongs to the King and summons the recovered person toward the Father.\\
St. Gregory of Nyssa, \emph{On Virginity}, 12 & The lamp becomes reason illumined for an inward search; the house is the person's interior life. & The royal image has been covered by the accumulated dirt of disordered life, not annihilated. & The soul searches its faculties, uncovers the image, and calls the ordered powers of the soul to rejoice.\\
\end{fourstable}

\noindent Gregory and Ambrose expound Luke 15 directly. Gregory of Nyssa reuses the parable in an ascetical account of self-knowledge. The three readings work at compatible levels and remain distinct: divine Wisdom, the Church, and the searching soul are not interchangeable subjects in the Fathers' arguments.

\begin{massbox}{Outside and inside}
The sheep prevents repentance from becoming mere introspection: the Shepherd comes from beyond the wanderer's resources. The drachma prevents lost from meaning worthless: the King's likeness remains the ground of the search. Together they describe rescue as pursuit, illumination, cleansing, and restored communion.
\end{massbox}

\sectionguard
\subsection*{Repentance and the enlargement of joy\enspace\properrefs{Gosp., Off., Comm.}}

\begin{threestable}{Movement}{Patristic precision}{Liturgical placement}
Mourning becomes changed life & Gregory defines repentance by joining grief over evil already done to refusal to repeat what must be grieved. & Luke's explicit sinner doing penance governs the Communion's joy; reception cannot be detached from conversion.\\
Memory becomes ardor & Gregory notes that some converts love more intensely because they remember the danger from which they were rescued. & The Offertory's trust and praise are the found person's gratitude, not a boast over those who did not visibly wander.\\
Private recovery becomes common feast & Augustine, \emph{Confessions} VIII.3--4, applies Luke 15 to Victorinus's conversion and asks why the whole Church rejoiced so publicly. A widely known return naturally enlarges the circle that receives the news. & Greater publicity explains the breadth of communal joy; it does not make a prominent convert more valuable before God than an obscure one.\\
Angelic joy corrects human comparison & Cyril and Gregory require the hearer to rejoice at the good accomplished in another. & The Communion makes the assembly sing the answer the murmurers withheld.\\
\end{threestable}

\begin{studybox}{Grace and response remain braided}
The seeker acts first and until the lost is found. Yet Luke names repentance, Peter commands humility and resistance, and the altar asks for sanctification. Augustine's Victorinus episode makes one recovered life a good shared by the Church while expressly denying that fame makes that soul more valuable. Patristic praise of prevenient mercy does not make the recovered person inert; moral exhortation does not make the sinner his own shepherd.
\end{studybox}

\sectionguard
\subsection*{Entrusted care and sober resistance\enspace\properrefs{Ep., Grad., All.}}

First Peter 5:6--11 contains a tightly ordered sequence. Humility places the Christian under God's mighty hand; care is cast upon him because he cares; sobriety recognizes a real adversary; resistance is made in faith and in solidarity with the suffering brotherhood; God himself restores, strengthens, and establishes.

\begin{fourstable}{Apostolic act}{Its ground}{Error excluded}{Liturgical echo}
Humble yourselves & God's mighty hand will exalt in due time & Humility as self-erasure or despair & The Introit discloses poverty to a God asked to look and pity.\\
Cast all care upon him & He has care of you & Anxious control or irresponsible passivity & The Gradual sings Psalm 54: \latin{Iacta cogitatum tuum in Domino}.\\
Be sober and watch & The devil seeks someone to devour & Naive denial of spiritual danger & The solitary person of the Introit is gathered into a suffering brotherhood.\\
Resist, steadfast in faith & The same sufferings are borne throughout the brotherhood & Isolated heroism or surrender to fear & The Offertory turns the recovered person toward trust, seeking, praise, and the poor.\\
Endure in hope & The God of all grace will perfect and establish & Treating present affliction as the final word & The Postcommunion asks that holy things give life and prepare for everlasting mercy.\\
\end{fourstable}

\begin{sourcecard}{St. Augustine --- on Psalm 54:23}
In \emph{Exposition on Psalm 54} (Psalm 55 in modern numbering), 23, Augustine directly explains the Gradual's \latin{Iacta cogitatum tuum}. The burden is not handed to another fragile human being as though that person were the final support; the self and its care are entrusted to the Creator who made and sustains it. First Peter's adjacent command echoes the same verse; applying Augustine's exposition to it is a reading; Augustine does not comment on the Epistle.
\end{sourcecard}

The other appointed psalms add three points that the Gradual alone does not supply. Augustine's Psalm 24 exposition lets the solitary and poor speaker lift the soul without claiming self-sufficiency: shame is refused through trust, while affliction and sins remain confessed. On Psalm 7 he treats the just Judge's patience as preserving medicine and makes received righteousness, not autonomous merit, the ground of confidence. On Psalm 9 he joins knowledge of God's name to seeking, proclamation from Sion, divine remembrance of blood, and the poor person's cry. Thus Introit, Alleluia, and Offertory move from exposed need through patient judgment into praise that remembers other sufferers.

Anthony of Padua later preached the same Gospel and Epistle as one three-part concordance. Publicans approach, the Shepherd carries, and the woman lights the house; humility, vigilance, and God's final strengthening answer those movements from 1~Peter. His most distinctive move is to make the lamp's clay, wick, oil, and flame a searching discipline of frailty, penance, mercy, and divine love. Gregory and Ambrose expound Christ and Church more directly; Anthony shows how their received images became medieval examination of conscience without reducing the lost coin to private self-rescue.

\begin{sourcecard}{Tertullian --- patience is pursuit}
\emph{On Patience} 12 sets the patient Shepherd against the impatience that would despise one lost sheep. This corrects a passive reading of the Alleluia's \latin{patiens}: divine patience seeks, bears, and restores, while Christian patience remains sober and steadfast under the same God.
\end{sourcecard}

\begin{massbox}{The predator is not the Sunday's deepest actor}
The devil seeks in order to consume; Christ seeks in order to carry home. Vigilance is urgent because danger is real, but hope is rational because the God of all grace has promised the last verbs: restore, strengthen, and establish.
\end{massbox}

\sectionguard
\subsection*{Under a ruler and guide: temporal goods without temporal captivity\enspace\properrefs{Coll., Off., Sec., Postcomm.}}

The Collect is not a generic request for protection. Its clauses distinguish the source of strength and holiness, the government of the journey, the goods traversed, and the end that must not be lost.

\begin{threestable}{Clause}{Textual force}{Formulary consequence}
\latin{sine quo nihil est validum, nihil sanctum} & Apart from God nothing possesses the strength or holiness named by the prayer. & The petition begins from dependence rather than from confidence in the traveller's unaided capacity.\\
\latin{multiplica super nos misericordiam tuam} & Mercy is asked to be multiplied over the worshippers. & Divine government is itself mercy before it becomes a rule for choosing among goods.\\
\latin{te rectore, te duce} & The repeated ablative absolute names God as ruler and guide. & Direction is both normative and active: God governs the end and conducts the passage.\\
\latin{per bona temporalia} & The prayer asks to pass through temporal \emph{goods}, not through things redescribed as evil. & The Offertory's praise and the Secret's gifts can belong to faithful passage; created goods are used without becoming the final good.\\
\latin{ut non amittamus aeterna} & The negative final clause fixes the danger: eternal goods may be lost. & The Postcommunion's everlasting mercy names the horizon toward which present reception must prepare the faithful.\\
\end{threestable}

\begin{studybox}{Neither contempt nor captivity}
The Collect calls temporal things good, excluding contempt for creation; because they are traversed under God's rule toward eternal goods, it also excludes captivity to them. Read within the one formulary, the later gifts and holy things continue this itinerary.
\end{studybox}

\sectionguard
\subsection*{Gifts looked upon, believers sanctified\enspace\properrefs{Int., Coll., Sec., Postcomm.}}

\begin{studybox}{The Secret's Latin and English witnesses}
The 1962 Latin runs from \latin{Respice, Domine} to its printed endpoint at \latin{in unitate}. Lefebvre's historical English of 1925 has its own abbreviated conclusion, ``Through our Lord.'' The two therefore end differently, each as printed.
\end{studybox}

\begin{threestable}{Latin hinge}{Grammatical force}{Place in the Sunday's movement}
\latin{Respice, Domine} & An imperative asking for God's attentive and favorable regard & The Introit's personal \latin{Respice in me} reaches the altar as the Church's \latin{Respice \ldots{} munera}. The sought sinner is incorporated into a supplicating offerer.\\
\latin{munera supplicantis Ecclesiae} & The gifts of the supplicating Church are the object of the divine gaze & The prayer is ecclesial, not a private transaction; the Church offers while still begging.\\
\latin{saluti credentium} & A dative of end or advantage: for the salvation of believers & The gifts' requested fruit is personal and corporate salvation, not visible success or possession.\\
\latin{perpetua sanctificatione} & An ablative naming the enduring sanctification with which reception is joined & Sanctification is divine action, not a quality the gifts or offerers manufacture independently.\\
\latin{sumenda concede} & The gerundive agrees with \latin{munera}: grant the gifts to be received & \latin{sumenda} anticipates the Postcommunion's \latin{sumpta}: what is petitioned before the canon is confessed after reception.\\
\end{threestable}

\begin{sourcecard}{St. Irenaeus --- oblation and reception}
\emph{Against Heresies} IV.18.1--5 teaches that the Church offers from God's own creation, that acceptance glorifies the offerer rather than supplying a divine need, and that Eucharistic bread receives a heavenly reality. Bodies nourished by it are therefore not destined for annihilation but for resurrection. Irenaeus does not comment on this Secret; his doctrine frames its movement from ecclesial gifts to sanctification and salvation.
\end{sourcecard}

\begin{sourcecard}{St. Cyril of Jerusalem --- sanctified gifts and holy reception}
\emph{Mystagogical Catechesis} 5 (Catechetical Lecture 23), especially 7 and 19--21, describes the invocation of the Holy Spirit upon the gifts, the proclamation of holy things for the holy, and the communicant's reverent reception. Cyril's mystagogy belongs to a different local rite and does not address the Roman prayer. It illuminates why sanctification and reception belong together without making reception automatic or morally indifferent.
\end{sourcecard}

\begin{studybox}{Four verbal threads through the formulary}
\begin{itemize}
\item \latin{Respice in me} becomes \latin{Respice \ldots{} munera}: the person under mercy enters the offering Church.
\item The Collect's \latin{sperantium} and the Secret's \latin{credentium} hold hope and faith within God's protection.
\item \latin{perpetua sanctificatione} opens into the Postcommunion's \latin{Sancta tua}: sanctification is received in holy things, not reduced to an inward idea.
\item \latin{sumenda} becomes \latin{sumpta}: petition before reception becomes thanksgiving and a new petition for vivification afterward.
\end{itemize}
\end{studybox}

\sectionguard
\subsection*{Sacramental completion: joy, life, and everlasting mercy\enspace\properrefs{Int., Gosp., Off., Sec., Comm., Postcomm.}}

The Communion repeats Luke 15:10 rather than choosing an unrelated sacramental verse. The assembly receives while singing heaven's joy over repentance. The Postcommunion then refuses to let joy remain a momentary emotion: the holy things received are asked to vivify, purify, and prepare the faithful for mercy without end.

\begin{threestable}{Moment}{Change of subject}{Sacramental consequence}
Introit & The isolated sufferer asks God to look upon him. & Need is brought into truthful presence before God.\\
Gospel & Christ reveals the Shepherd and Wisdom who seek until they find. & Divine initiative judges both despair and contempt.\\
Offertory & Those who know God's name trust, seek, praise, and remember the poor. & The recovered person acquires an ecclesial voice and concern for other sufferers.\\
Secret & The supplicating Church asks God to look upon her gifts. & Personal rescue becomes common oblation ordered to believers' salvation.\\
Communion & The Church sings the angels' joy over one sinner repenting. & The murmurer's gaze is replaced by a received, heavenly pattern of desire.\\
Postcommunion & The Church speaks as those who have received holy things. & The gift is asked to give life now and prepare the purified for everlasting mercy.\\
\end{threestable}
